\documentclass[11pt]{scrartcl}
\usepackage[sexy]{evan} %BLBLBLBLBLBLBLBLBLBLBLB EVANCHEN
\usepackage[T1]{fontenc}
\usepackage[utf8]{inputenc}
\usepackage{graphicx}
\usepackage{amsmath}
\usepackage{tikz}
\usepackage{tikz-cd}
\usepackage{asymptote}
\usepackage{amsthm}
\usepackage{amssymb}
\usepackage{gensymb}
\usepackage{amsfonts}
\usepackage[english,italian]{babel}
\usepackage{quoting}
\quotingsetup{font=big}
\usepackage{booktabs}
\usepackage{caption}
\usepackage{lmodern}
\usepackage{faktor}
\usepackage{bbm} %oer fare gli insiemi indicatori
\usepackage{leftindex} %per gli ideali e gli indici a sinistra
\usepackage{hyperref}
\usepackage{epigraph} 
\usepackage{pgfplots}
\usepackage[export]{adjustbox}

%Pacchetti per i codici
\usepackage{listings}
\usepackage{xcolor}
\usepackage{algpseudocode}

%Analisi e probabilità
%----------------------------------
\newcommand{\dif}{\text{d}}
%%% the pure symbol, see https://tex.stackexchange.com/a/718194/4427
\newcommand{\cind}{\perp\mspace{-9mu}\perp}
%%% define \indep
\NewDocumentCommand{\indep}{e{_}}{%
	\mathrel{%
		\IfNoValueTF{#1}{% no subscript
			\cind
		}{% subscript
			\mathchoice{\mathop{\cind}\limits_{#1}}{\cind_{#1}}{\cind_{#1}}{\cind_{#1}}%
		}%
	}%
}
%-----------------------------------

%Algebra lineare e geometria
%----------------------------------
\newcommand{\tr}[1]{\text{tr}\!\left(#1\right)}
\renewcommand{\rank}[1]{\text{rank}\!\left(#1\right)}
\newcommand{\Immagine}[1]{\text{Im}\left(#1\right)}
\renewcommand{\ker}[1]{\text{Ker}\left(#1\right)}
\renewcommand{\AA}{\mathbb{A}}
\newcommand{\LL}{\mathbb{L}}
\newcommand{\SO}{\text{SO}}
%----------------------------------

%Fisica
%----------------------------------
\newcommand{\UDiMisura}[1]{\,\text{#1}}
\newcommand{\ihat}{\hat{\imath}}
\newcommand{\jhat}{\hat{\jmath}}
\newcommand{\khat}{\hat{k}}
\newcommand{\posit}{\vec{r}}
\newcommand{\veloc}{\dot{\vec{r}}}
\newcommand{\accel}{\ddot{\vec{r}}}
%----------------------------------

%Informatica
%----------------------------------
\definecolor{codegreen}{rgb}{0,0.6,0}
\definecolor{codegray}{rgb}{0.5,0.5,0.5}
\definecolor{codepurple}{rgb}{0.58,0,0.82}
\definecolor{backcolour}{rgb}{0.95,0.95,0.92}

\lstdefinestyle{overleafwiki}{
	backgroundcolor=\color{backcolour},   
	commentstyle=\color{codegreen},
	keywordstyle=\color{magenta},
	numberstyle=\tiny\color{codegray},
	stringstyle=\color{codepurple},
	basicstyle=\ttfamily\footnotesize,
	breakatwhitespace=false,         
	breaklines=true,                 
	captionpos=b,                    
	keepspaces=true,                 
	numbers=left,                    
	numbersep=5pt,                  
	showspaces=false,                
	showstringspaces=false,
	showtabs=false,                  
	tabsize=2
}

\lstset{style=overleafwiki}
\newcommand{\bigo}[1]{\mathcal{O}\!\left(#1\right)}
\newcommand{\bigomega}[1]{\Omega\!\left(#1\right)}
\newcommand{\bigtheta}[1]{\Theta\!\left(#1\right)}
%----------------------------------

%Insiemistica e Algebra
%----------------------------------
\newcommand{\eset}{\varnothing}
\newcommand{\parti}[1]{\mathcal{P}\!\left(#1\right)}
\newcommand{\inj}{\hookrightarrow}
\newcommand{\surj}{\twoheadrightarrow}
\newcommand{\bij}{\hookrightarrow\mathrel{\mspace{-15mu}}\rightarrow}
\newcommand{\ordine}[2]{\text{ord}_{#1}\left(#2\right)}
\newcommand{\normal}{\mathrel{\unlhd}}
\newcommand{\MCD}{\text{MCD}}
\renewcommand{\mcm}{\text{mcm}}
%----------------------------------


\begin{document}
	\title{Appunti di Teoria della misura}
	\subtitle{Elio Marconi - SGSS}
	\date{Settembre - Dicembre 2026}
	
	
	\maketitle
	
	\tableofcontents
	
	\newpage
	
	\section{Preambolo: misure in $\RR^n$}
	
	Questa sezione è stata scritta leggendo "Real Analisys" di E.M. Stein e R. Shakarchi. L'idea è, prima di attaccare la teoria della misura nella sua piena astrazione, di ottenere un po' di intuizione nello spazio euclideo $\RR^n$ che siamo in grado di visualizzare. 
	
	\subsection{La misura esterna}
	
	In $\RR^n$ la prima cosa originale che possiamo definire è 
	\begin{definition}[Cubi e Rettangoli]
		In $\RR^n$ dati gli intervalli
		$$I_n=\left[a_n,b_n\right]$$
		definiamo il \textit{rettangolo chiuso} definito da questi intervalli come
		$$R:= I_1\times I_2 \times \dots \times I_n$$
		analogamente possiamo definire il \textit{rettangolo aperto} sostituendo agli intervalli chiusi degli intervalli aperti. Il \textit{volume} di $R$ è definito come
		$$\abs{R}:= \prod_{i=1}^{n}(b_n-a_n)$$
		un rettangolo tale che $b_i-a_i=b_j-a_j$ per ogni coppia di indici $i,j$ si dice un \textit{cubo}.
	\end{definition}  
	Con questo possiamo definire una nozione di misura esterna, che non esiste in spazi arbitrari:
	\begin{definition}[Misura esterna]
		Dato un qualunque $E\subseteq\RR^n$ definiamo la sua \textit{misura esterna} $m_*:\parti{\RR^n}\to [0,\infty]$ come:
		$$m_*(E)=\inf\sum_{i=1}^{\infty}\abs{Q_i}$$
		dove l'inf è preso su tutti i possibili ricoprimenti $Q_i$ di $E$ con $Q_i$ cubi chiusi, ovvero abbiamo che 
		$$E\subseteq \bigcup_{i=1}^\infty Q_i$$
		con ciascuno dei $Q_i$ un cubo chiuso.
	\end{definition}
	
	
\end{document}
